% BHU PYQ -- Manually transcribed & error-corrected
\documentclass[11pt,a4paper]{article}

\usepackage[a4paper,top=2.5cm,bottom=2.5cm,left=2.8cm,right=2.8cm,headheight=14pt]{geometry}
\usepackage{amsmath,amssymb}
\usepackage[T1]{fontenc}
\usepackage[utf8]{inputenc}
\usepackage{lmodern}
\usepackage{microtype}
\usepackage{fancyhdr}
\usepackage{enumitem}
\usepackage{hyperref}
\usepackage{lastpage}
\usepackage{parskip}

\hypersetup{colorlinks=true,urlcolor=black,linkcolor=black,
  pdftitle={CHB-501: Analytical Chemistry-I -- BHU PYQ},pdfauthor={Banaras Hindu University}}

\pagestyle{fancy}\fancyhf{}
\renewcommand{\headrulewidth}{0.4pt}\renewcommand{\footrulewidth}{0.4pt}
\fancyhead[L]{\small   \textit{Banaras Hindu University}}
\fancyhead[R]{\small   \textit{CHB-501: Analytical Chemistry-I}}
\fancyfoot[C]{\small Page  \thepage\ of \pageref{LastPage}}


\newlist{parts}{enumerate}{2}
\setlist[parts,1]{label=(\alph*),leftmargin=*,itemsep=5pt,topsep=3pt}
\setlist[parts,2]{label=(\roman*),leftmargin=*,itemsep=3pt,topsep=2pt}

\newcommand{\pts}[1]{\hfill   \textit{[#1]}}


\begin{document}


\begin{center}
  {\large\bfseries BANARAS HINDU UNIVERSITY}\\[2pt]
  {\normalsize B.Sc. (Hons.) Semester V Examination, 2022-23}\\[6pt]
  \rule{0.8\linewidth}{0.5pt}\\[6pt]
  {\large\bfseries Paper No. CHB-501: Analytical Chemistry-I}\\[4pt]
  {\normalsize Time: 3 Hours\hfill Maximum Marks: 70}
\end{center}
\rule{\linewidth}{0.4pt}
\small



\noindent   \textit{Instructions: Answer five questions in total, including Question~1 which is compulsory.}

\normalsize
\bigskip




\noindent   \textbf{Question 1.} Answer all parts of the following: \pts{5 $ \times$ 2 = 10}

\begin{parts}
  \item What is the difference between specific and selective organic precipitants?
  \item What is the role of oxine in gravimetric analysis?
  \item The water of crystallization in $ \text{CuSO}_4\cdot5 \text{H}_2 \text{O}$ could be determined by thermogravimetric analysis. Explain.
  \item Complete the following Malaprade reaction:
    \[  \text{CH}_2 \text{OH-CH}_2 \text{OH} +  \text{IO}_4^-  o ? \]
  \item Brief the key features of Good Laboratory Practices (GLP) and Good Manufacturing Practices (GMP).
\end{parts}

\medskip\rule{\linewidth}{0.2pt}




\noindent   \textbf{Question 2.}

\begin{parts}
  \item Discuss various steps involved in gravimetric analysis. What is the von Weimarn ratio? Define various terms used in it. \pts{8}
  \item Write the use of any two of the following organic precipitants in chemical analysis: (i) Dithizone, (ii) Oxine, and (iii) Cupferron. \pts{6}
\end{parts}

\medskip\rule{\linewidth}{0.2pt}




\noindent   \textbf{Question 3.}

\begin{parts}
  \item What is the isomorphous replacement in the development of precipitates? Explain. \pts{4}
  \item How can you minimize the magnitude of determinate errors in chemical analysis? \pts{4}
  \item Differentiate between accuracy and precision. Calculate the absolute and relative precision parameters of the following set of results obtained during the determination of the chromium content of industrial wastewater: $6.8, 7.7, 7.2, 7.3$, and $7.5\, \text{ppm}$. \pts{6}
\end{parts}

\medskip\rule{\linewidth}{0.2pt}




\noindent   \textbf{Question 4.}

\begin{parts}
  \item Discuss the chemistry and stoichiometry of the estimation of vicinal diols using periodic acid. How do you standardize this reagent? \pts{8}
  \item How can you estimate $\alpha$-hydroxyamines using Malaprade reagent? \pts{6}
\end{parts}

\medskip\rule{\linewidth}{0.2pt}




\noindent   \textbf{Question 5.} How do you estimate the following? Write reactions: \pts{6 + 4 + 4 = 14}

\begin{parts}
  \item Metal ions using potassium bromate as a redox titrant
  \item Nitrite using Chloramine-T
  \item Anhydrides using Karl Fischer reagent
\end{parts}

\vfill
\rule{\linewidth}{0.4pt}

\begin{center}\small
  \href{https://bhu-exam.vercel.app/}{Click here to visit our website}\\[4pt]
  \href{https://me-aryan.vercel.app/}{Click here to visit our contributor}\\[4pt]
  \href{https://quashan-me.vercel.app/}{Click here to visit our contributor}
\end{center}

\end{document}
